\section*{Temporal Clauses}

An explicit ``when P, Q'' construction uses \textit{tu-le-ngothr}, ``at the time,'' followed by a finite event description:

\begin{quote}
\texttt{tu-le-ngothr: P; Q.}
\end{quote}

This means ``At the time described by P, Q.'' The head combines existing \textit{tu-} (temporal case), \textit{le-} (abstract class), and \textit{ngothr} (time). The finite clause P describes the relevant time. No new lexical word is introduced.

\subsection*{Boundaries and scope}

\begin{itemize}
\item The colon opens the event description P.
\item The semicolon closes P and introduces the main clause Q.
\item Commas can join events inside P: \textit{tu-le-ngothr: P1, P2; Q}. Both events fall inside the temporal frame. Order alone does not explicitly mean ``and then.''
\item In speech, continuing intonation after the time expression and a pause with an intonation reset before Q mark the same scope. Repeat a subject when contextual recovery would be unclear.
\end{itemize}

These punctuation roles apply inside this construction. Outside an open temporal frame, ordinary asyndetic coordination still applies. These examples establish one frame; deeper nested temporal frames need a separately specified grouping convention.

\subsection*{Finite clauses and agreement}

P and Q each use ordinary subject, case, and auxiliary rules. Their subjects and person prefixes may differ. Every auxiliary retains person marking, and any omitted subject must be recoverable from context.

\begin{itemize}
\item \textit{tu-le-ngothr: na ifa fethr; wa ofa for.}\\
``When I speak, you write.''
\item \textit{tu-le-ngothr: i-la-melhen ashe lheth; i-la-mar afe for.}\\
``When the child was sleeping, the mother wrote.''
\end{itemize}

Each clause contributes its own tense and aspect; agreement does not require matching tense or aspect across P and Q. The temporal interpretation must remain coherent. The frame alone does not assert completion of P, causation, or conditionality; aspect and context supply those distinctions.

\subsection*{Other constructions and open questions}

\textit{Tu-lhan} remains the interrogative ``when?'' Bare clause juxtaposition can still receive a temporal reading from context. This explicit frame licenses a finite event description after its temporal head; general relative clauses, other noun-clause complements, and tense interpretation inside nested auxiliary chains remain separate questions.

In the working sentence, pondering and writing belong inside P; inheriting life, taking root, and becoming impossible to ignore belong in Q. ``Then,'' ``their own,'' ``take root,'' and ``ignore'' remain to be settled.
